%!TEX TS-program = xelatex
%!TEX encoding = UTF-8 Unicode
\documentclass[11pt,a4paper]{awesome-cv}
\usepackage[hidelinks]{hyperref}

% --- ATS / PDF metadata ---
\hypersetup{
  pdftitle={Juan Manuel Young Hoyos — Cover Letter — Principal Software Engineer},
  pdfauthor={Juan Manuel Young Hoyos},
  pdfsubject={Cover Letter — Principal Software Engineer — Autodesk},
  pdfkeywords={TypeScript, serverless, AWS, large-scale systems, platform, design automation, AI-enabled workflows, architecture, technical leadership, Autodesk, Forma}
}

\geometry{left=1.4cm, top=.8cm, right=1.4cm, bottom=1.8cm, footskip=.5cm}

% Emit ActualText so PDF text extraction yields ASCII hyphens (ATS keyword matching)
\XeTeXgenerateactualtext=1

\colorlet{awesome}{awesome-red}
\setbool{acvSectionColorHighlight}{true}
\renewcommand{\acvHeaderSocialSep}{\quad\textbar\quad}

\definecolor{darktext}{HTML}{1A1A1A}
\colorlet{text}{darktext}
\renewcommand*{\lettertextstyle}{\fontsize{10pt}{1.45em}\bodyfontlight\upshape\color{darktext}}
\renewcommand*{\lettersectionstyle}[1]{{\fontsize{14pt}{1em}\bodyfont\bfseries\color{darktext}#1}}

% --- Header info ---
\name{Juan Manuel}{Young Hoyos}
\position{Principal Software Engineer — TypeScript · Cloud \& Large-Scale Systems · Design Automation}
\address{Medellín, Colombia · Remote · Open to relocating to the Netherlands}
\mobile{+57 1234567890}
\email{juanmanuel12.13jmyh81@gmail.com}
\homepage{jmyounghoyos.com}
\github{Youngermaster}
\linkedin{juan-manuel-young-hoyos}

% --- Letter meta ---
\recipient
  {Forma Design — Hiring Team}
  {Autodesk}
\letterdate{\today}
\lettertitle{Application for Principal Software Engineer — Forma Design}
\letteropening{Dear Autodesk Hiring Team,}
\letterclosing{Sincerely,}
\letterenclosure[Attached]{Resume — Juan Manuel Young Hoyos.pdf}

\begin{document}
\makecvheader[R]
\makelettertitle

\begin{cvletter}

I am applying for the \textbf{Principal Software Engineer} role on the Forma Design team. The part of your description that made me want to write is turning design automations into \textbf{reusable, configurable building blocks} that users can apply without writing code or understanding the implementation. That is exactly the problem I am working on today, in a different domain.

As \textbf{Technical Lead at Okorum Technologies} I set the technical direction for a national-scale interoperability platform and delivered it ahead of a hard government deadline. I architected it as a \textbf{Turborepo monorepo of NestJS microservices} coordinating staged processing over \textbf{message queues}, engineered fault tolerance and resumability so work survives an unreliable upstream system, and stood up \textbf{CI/CD on AWS} with full observability. I stayed hands-on in the code the whole way through --- which is the part of this role I care most about keeping.

I have also driven architecture \textbf{beyond a single team}: as consulting architect through Grisú I defined system architecture for client organisations that other engineers then built against, and at Design Systems Inno I built the shared internal API platform --- an asynchronous, queue-backed design where clients track long-running jobs by ticket rather than blocking --- used across several product teams.

On \textbf{AI-enabled workflows}, I ship features to \textbf{Unicorn AI}, an internal AI-agent platform whose configurable per-role agents automate the work of developers, designers, accountants and lawyers. It is the same shape as your automations: reusable building blocks, adapted per context, with the complexity hidden from the user. I also use \textbf{Claude Code} and LLMs daily in my own engineering, while reviewing and validating everything they produce.

I lead through influence and example rather than authority --- mentoring is something I have done deliberately for over a decade, from tutoring 100+ students to leading engineering and QA teams. I would welcome the chance to bring that, and a genuine interest in sustainable design tooling, to Forma Design. Thank you for your consideration.

\end{cvletter}

\makeletterclosing
\end{document}
